\documentclass[10pt]{article}

\usepackage[margin=1in]{geometry}
\usepackage[T1]{fontenc}
\usepackage[utf8]{inputenc}
\usepackage{array}
\usepackage{booktabs}
\usepackage{enumitem}
\usepackage{hyperref}
\usepackage{xcolor}
\usepackage{longtable}

\hypersetup{colorlinks=true,linkcolor=blue,urlcolor=blue}
\setlist{nosep,leftmargin=*}
\newcommand{\code}[1]{\texttt{#1}}

\title{Artifact: Chronos: Scalable and Cost-Effective Cloud-based RAN Emulation System}
\author{}
\date{}

\begin{document}
\maketitle

\section{Artifact overview}

This artifact accompanies the paper \emph{Chronos: Scalable and Cost-Effective
Cloud-based RAN Emulation System}.  Its purpose is to make the implementation available for inspection,
provide the scripts used to deploy Chronos, provide the paper's plot outputs and
their underlying data, and demonstrate that the Chronos code is functional on a
Powder/Emulab deployment.

The evaluation combines two complementary forms of evidence.  First, the
artifact supplies the data and scripts used to generate the plots in the
paper.  Second, the hosted Powder deployment gives evaluators a hands-on
view of Chronos operating end to end: its components start, coordinate RAN
slots, attach a UE, and carry data-plane traffic.

A video tutorial covering the complete artifact and Powder evaluation workflow
will be available in the
\href{https://github.com/netsys-edinburgh/EmuRAN-Draft-Artifact}{GitHub
repository}.

\subsection{Badge scope}

The artifact is prepared to support all three requested badges:

\begin{itemize}
  \item \textbf{Artifacts Available:} the source, deployment automation,
  experiment data, and plotting scripts are publicly available in the GitHub
  repository.
  \item \textbf{Artifacts Evaluated--Functional:} evaluators connect to the
  hosted Powder deployment and verify Chronos's end-to-end operation.
  \item \textbf{Results Reproduced:} evaluators run the plotting notebook and
  reproduce the key scalability result on Powder by measuring dilation at
  increasing gNB/UE scale.
\end{itemize}

\section{Artifact components}

The artifact has four main components.

\subsection{Source code}

The \code{code/} directory contains the implementation of Chronos:

\begin{longtable}{>{\raggedright\arraybackslash}p{0.27\linewidth}>{\raggedright\arraybackslash}p{0.66\linewidth}}
\toprule
Path & Purpose \\
\midrule
\endhead
\code{code/supervisor/} & Patched Linux 5.15/KVM source.  The virtualized TSC allows guest time to be advanced or paused independently of wall-clock time. \\
\code{code/slot\_monitor/} & The \code{custom\_tsc} kernel module, shared-memory support, and slot-monitor utilities used to control virtual time. \\
\code{code/middlebox/} & Modified multi-UE nFAPI proxy.  It forwards traffic, tracks local slot completion, and communicates with the global slot coordinator. \\
\code{code/gnb/} & Modified OpenAirInterface gNB source, operated in nFAPI VNF/emulated-L1 mode. \\
\code{code/ue/} & Modified OpenAirInterface UE source, operated in standalone nFAPI PNF/emulated-L1 mode. \\
\bottomrule
\end{longtable}

Together these components implement PHY bypass, virtual time, slot-level
synchronization, and distributed control/data-plane forwarding.  The gNB, UE,
and supervisor directories are complete source trees rather than small patch
sets so that the evaluated implementation is preserved in full.

\subsection{Deployment scripts}

The \code{deployment/} directory contains the automation used to prepare and
run Chronos:

\begin{itemize}
  \item \code{deployment/provisioning/profile.py} is an Emulab/Powder profile
  that allocates controller/compute nodes, proxy nodes, and a global slot
  coordinator.
  \item \code{deployment/provisioning/scripts/} configures Ubuntu hosts,
  builds and installs the patched kernel, creates nested KVM guests, installs
  the virtual-time module and services, configures networking, and constructs
  the Kubernetes cluster.
  \item \code{deployment/helm/} contains Helm charts for the mobile core,
  global slot coordinator, proxy, gNB, UE, and persistent-UE components.
  \item \code{deployment/helm/values.yaml} controls topology size, image
  names, addressing, startup behavior, and the UE data-plane workload.
  \item Operational scripts under \code{deployment/helm/} start UEs, inspect
  or restart components, and collect logs.
\end{itemize}

These scripts perform privileged operations and modify kernels, networking,
storage, VMs, and cluster state.  They must only be used on disposable
Powder/Emulab experiment nodes.

\subsection{Plot outputs}

The \code{results/} directory has two subdirectories:

\begin{itemize}
  \item \code{results/data/} contains raw logs, packet captures, CSV/JSON
  intermediates, and analysis/plotting scripts from the paper experiments.
  \item \code{results/output/} contains the paper's prepared PDF plot outputs.
\end{itemize}

The top-level \code{plot.ipynb} provides a convenient way to inspect the plot
data, plotting scripts, and PDF outputs.  Evaluators may browse the prepared
outputs directly and trace each result back to its associated analysis code
and recorded measurements.

\subsection{Running the code and deployment scripts on Powder/Emulab}

The authors run a functional deployment created with the supplied
Powder/Emulab profile and deployment automation.  Evaluators receive SSH
access to this running deployment and inspect its operation.  The goal is to
demonstrate that the implementation and deployment scripts work together.  The
preserved experiment data and outputs complement this live demonstration
with evidence from the paper's full-scale evaluations.

\section{Running the plotting notebook}

The top-level \code{plot.ipynb} runs the plotting script associated with each
paper figure.  Each code cell handles one figure or panel, reads its inputs
from \code{results/data/}, writes a newly generated PDF to
\code{results/output/}, and displays that PDF inline.  The notebook uses paths
relative to the artifact root and therefore does not require any
machine-specific path changes.

\subsection{Obtaining the artifact}

Clone the artifact repository from GitHub and enter its root directory:

\begin{verbatim}
git clone https://github.com/netsys-edinburgh/EmuRAN-Draft-Artifact.git
cd EmuRAN-Draft-Artifact
\end{verbatim}

The repository is available at
\href{https://github.com/netsys-edinburgh/EmuRAN-Draft-Artifact}{github.com/netsys-edinburgh/EmuRAN-Draft-Artifact}.
All installation and notebook commands below should be run from this directory.

\subsection{Software dependencies}

The notebook runs on Linux and macOS with Python 3.10 or newer.  On Linux,
install Python 3 and its \code{venv} support using the distribution package
manager.  On macOS, use a current Python installation from
\href{https://www.python.org/}{python.org} or Homebrew.  The notebook and
plotting scripts require the following Python packages:

\begin{itemize}
  \item \code{jupyter}, \code{IPython}, and \code{ipykernel} for the notebook
  interface, inline display, and the dedicated artifact kernel;
  \item \code{numpy} for numerical processing;
  \item \code{pandas} for CSV and tabular data processing;
  \item \code{matplotlib} and \code{seaborn} for plotting; and
  \item \code{pymupdf} for rendering each generated PDF inside the notebook.
\end{itemize}

No GPU, LaTeX installation, or privileged access is needed for this notebook
on either platform.
Package requirements are listed in \code{requirement.txt}.  The supplied
installer verifies Python 3.10 or newer, creates \code{.venv}, and installs all
packages into that environment.  From the artifact root, run:

\begin{verbatim}
./install_dependencies.sh
\end{verbatim}

To use a different Python executable or virtual-environment location, set
\code{PYTHON\_BIN} or \code{VENV\_DIR}, respectively, before running the
installer.

Installing these packages requires Internet access once.  Subsequent notebook
runs use the data already included in the artifact.

\subsection{Launching and running the notebook}

Start Jupyter from the artifact root, which is the directory containing
\code{plot.ipynb} and \code{results/}:

\begin{verbatim}
source .venv/bin/activate
jupyter notebook plot.ipynb
\end{verbatim}

In the browser, select the \code{Python (Chronos Artifact)} kernel registered
by the installer.  This ensures that PyMuPDF and the other plotting
dependencies are available.  Run an individual cell to generate one figure,
or select \emph{Run All} to execute
the complete notebook.  The plotting programs use Matplotlib's
non-interactive \code{Agg} backend, so output appears inline rather than in a
separate desktop window.

Before each plot, the cell displays the paper figure number and the plotting
script being executed.  After the script completes, the cell renders the new
PDF and prints its path below the image.  A complete run produces 30 PDF files
under \code{results/output/}.  Re-running a cell replaces that cell's named
output while leaving the source data unchanged.

The notebook can also be executed non-interactively:

\begin{verbatim}
source .venv/bin/activate
jupyter nbconvert --to notebook --execute plot.ipynb \
  --output plot.executed.ipynb --ExecutePreprocessor.timeout=300
find results/output -maxdepth 1 -name '*.pdf' | wc -l
\end{verbatim}

The expected PDF count is \code{30}.  If a cell reports that it cannot find
\code{results/}, stop Jupyter and launch it again from the artifact root.  If
a plotting script fails, the cell reports the script's exit code and relevant
error message.

\section{Supported functionality}

The hosted demonstration checks the following properties:

\begin{enumerate}
  \item the Powder/Emulab profile allocates the required node roles;
  \item the provisioning scripts install the patched host/guest environment
  and create a working Kubernetes cluster;
  \item the core, global slot coordinator, proxy, gNB, and UE containers start;
  \item the gNB connects to the core and the UE completes registration;
  \item the local and global coordinators advance the emulation through RAN
  slots; and
  \item the attached UE obtains a working data plane and exchanges test
  traffic.
\end{enumerate}

Passing these checks demonstrates that the artifact code is functional.  The
plotting notebook regenerates the paper figures from the supplied data, and the
Powder scale sweep below reproduces the key relationship between emulation
scale, available compute, and the dilation factor.

\section{Access and evaluation window}

The authors provide a running Chronos deployment on Powder for the artifact
evaluation period from \textbf{September 1 through September 15}.  The authors
handle the Powder reservation, profile instantiation, kernel build, Kubernetes
cluster, and repository dependencies so evaluators can begin directly with
the functional checks.

To obtain access, each evaluator should post their \textbf{public SSH key} in
a comment on the paper's HotCRP artifact-evaluation page.  The authors will add
the key to the Powder deployment and reply through HotCRP with the SSH command,
host name, user name, and any required jump-host instructions.  Only a public
key (for example, the contents of \code{id\_ed25519.pub}) should be posted.
Evaluators must never send a private key, passphrase, password, or access token.

Access is available only during the stated September 1--15 window.  Requests
should be posted early enough to allow the authors to install the key and reply
within the evaluation period.  The evaluator needs an SSH client and the
private key corresponding to the submitted public key.

A video tutorial demonstrating the complete Powder evaluation workflow will
be available in the artifact's
\href{https://github.com/netsys-edinburgh/EmuRAN-Draft-Artifact}{GitHub
repository}.

\section{Evaluation on Powder/Emulab}

\subsection{Step 1: request and test SSH access}

Post the evaluator's public SSH key in a HotCRP comment.  After the authors
reply with connection details, use the supplied command to log in.  For
example, the command will have the following general form (the actual values
come from HotCRP):

\begin{verbatim}
ssh -i <private-key> <user>@<host>
\end{verbatim}

Do not change SSH authorization files or share the connection details outside
the artifact-evaluation process.  If login fails, report the exact SSH error
through HotCRP; do not post the private key.

\subsection{Step 2: connect to the controller}

The SSH command supplied through HotCRP connects to \code{node0}.  Node 0 is
the entry node for the experiment.  From \code{node0}, connect to the
controller and enter its Chronos deployment directory:

\begin{verbatim}
ssh ubuntu@10.2.1.2
cd chronos-auto-deploy
\end{verbatim}

\subsection{Step 3: configure and deploy one scale}

This experiment reproduces a key paper result: as the number of emulated
gNB/UE pairs increases on a fixed compute allocation, the work required per
RAN slot eventually exceeds the available real-time compute budget.  Chronos
responds by increasing its dilation factor, allowing more wall-clock time per
unit of virtual time while preserving coordinated RAN execution.

Evaluate the following gNB/UE pair counts in order:

\begin{verbatim}
1, 10, 50, 100, 200, 300
\end{verbatim}

For the first scale, edit \code{values.yaml}.  Set \code{numberGNB} and
\code{numberUE} to the same value.  Deploy that configuration with:

\begin{verbatim}
./run_experiment.sh
\end{verbatim}

When the interactive interface appears, press \code{p} to display pod status.
Wait until all expected pods are running.

\subsection{Step 4: verify connections and observe dilation}

After all pods are running, execute the supplied core-log helper:

\begin{verbatim}
./core-logs.sh
\end{verbatim}

Confirm from these logs that all configured gNBs and UEs are connected.  In
the interactive interface, enter \code{l globalsc} to display the global
coordinator log.  Observe and record the dilation value for the current scale.
A value near 1 means that the emulation is keeping pace with real time; a value
above 1 means Chronos is allocating additional wall-clock time to complete the
virtual slots.

\subsection{Step 5: tear down and repeat}

Remove the current deployment before changing scale:

\begin{verbatim}
./teardown-experiment.sh
\end{verbatim}

Wait for teardown to complete.  Then edit \code{values.yaml}, set both
\code{numberGNB} and \code{numberUE} to the next value in the sequence, and
repeat Steps 3--5.  Do not start a new deployment until the previous one has
been removed.

The result is reproduced when the measurements show the same key trend as the
paper: dilation remains close to real time while the fixed compute allocation
has capacity, then increases as gNB/UE scale creates compute pressure.  Exact
values can vary with Powder hardware and load; the reproduced result is the
increase in dilation that enables Chronos to cope with insufficient compute at
higher scale while continuing coordinated slot execution.  After recording
the value for scale 300, run \code{./teardown-experiment.sh} once more and log
out.  Report failures, missing permissions, or questions through HotCRP.

\section{Limitations and safety notes}

\begin{itemize}
  \item Cloud/testbed scheduling and network conditions may affect startup
  time and observed performance; evaluators can record these conditions when
  interpreting a live run.
  \item The deployment is shared during the evaluation window.  Evaluators
  should avoid disruptive or privileged commands and report unhealthy
  components through HotCRP.
\end{itemize}

\section{Summary of evaluation}

The evaluator should inspect the supplied code, deployment scripts, plot
outputs, and recorded data, then connect to the authors' running Powder
deployment during September 1--15.  Access requires only a public SSH key sent
through HotCRP comments.  The expected outcome is evidence that the supplied
Chronos implementation executes end to end with working slot coordination,
registration, and data-plane traffic.  Evaluators also regenerate the paper
plots from the supplied data and scripts and reproduce the key Powder result:
as gNB/UE scale increases on fixed compute, Chronos increases its dilation
factor to provide the wall-clock time required for coordinated slot execution.

\end{document}
